% Compile with the Dagstuhl OASIcs class:
%   \documentclass[a4paper,UKenglish,cleveref,autoref]{OASIcs}
% Place oasics.cls in the same directory (download from Dagstuhl).
\documentclass[11pt]{article}
\usepackage[utf8]{inputenc}
\usepackage{hyperref}
\usepackage{booktabs}
\title{Fused CPU-GPU Computing on RISC-V: A Unified RVV Architecture for Edge and Embedded Platforms}
\author{George Karani\\FrontForum Focus\\\texttt{georgekarani44@gmail.com}}
\date{PARMA-DITAM 2027}
\begin{document}
\maketitle
\begin{abstract}
Discrete GPUs provide high throughput for graphics and machine-learning workloads, but they impose separate instruction sets, memory spaces, and host-device data movement that are increasingly costly for energy-constrained edge and embedded platforms. This paper presents a fused CPU-GPU architecture built on RISC-V and its Vector (RVV) extension, in which general-purpose RISC-V cores equipped with wide vector units replace a discrete graphics processor. By unifying the instruction set, register-file conventions, and memory space, the design removes host-device transfers, simplifies the software stack, and enables fine-grained, energy-aware scheduling across heterogeneous workloads. We describe the microarchitecture, including a Tomasulo-style reorder buffer, per-thread L1 caches, sparsity-aware execution, and 128-bit packed INT8/INT4 vector units, and report on a Verilog prototype (tiny-gpu) targeting an AMD Xilinx Artix-7 FPGA carrier board. Preliminary simulation results demonstrate correct execution of matrix addition, matrix multiplication, and vector DP4A kernels with configurable thread counts. We discuss trade-offs, current limitations, and a roadmap toward cache coherence, branch divergence, and IEEE-754 FP32 support.
\end{abstract}
\noindent\textbf{Keywords:} RISC-V; RVV; unified CPU-GPU; heterogeneous computing; vector SIMD; open-source hardware; edge AI; energy efficiency.

\section{Introduction}
See paper.md for full text.

\end{document}
